% !TEX root = ../main.tex
% ============================================================
% Computational setup chapter
% Revised to match the frequency-transfer and results chapters.
% Focus: discrete systems, datasets, models, solvers, and
% residual-compatible diagnostics.
% ============================================================

\chapter{Computational setup}
\label{chapter:computational-setup}

\providecommand{\chsetupincludegraphics}[2]{%
    \IfFileExists{#2}{%
        \includegraphics[width=#1]{#2}%
    }{%
        \fbox{%
            \parbox[c][0.22\textheight][c]{0.86\linewidth}{%
                \centering
                \textbf{Figure placeholder}\\[0.6em]
                \texttt{\detokenize{#2}}%
            }%
        }%
    }%
}

The previous chapter introduced frequency transfer as a physically plausible
way to use a low-frequency Helmholtz solution in a high-frequency solve. It
also made the main solver-facing distinction: a learned field predictor is not
automatically a learned preconditioner. A Krylov method is controlled by
residuals, so every learned component must eventually be tested by the
discrete operator that the solver applies.

This chapter specifies the computational setup used to make that distinction
measurable. The purpose is not to report the experimental outcomes; those are
given in Chapter~\ref{chapter:results}. Instead, this chapter defines the
discrete Helmholtz systems, data conventions, neural transfer models,
preconditioned FGMRES protocol, and diagnostic quantities used later.

The experiments use three settings. The one-dimensional Dirichlet problem is a
controlled spectral microscope: its eigenvectors are known analytically, so
field-error components can be compared directly with residual components. The
one-dimensional PML problem is an intermediate test: it keeps the problem
small but introduces the absorbing-layer compatibility issue. The
two-dimensional finite-difference/PML problem is the solver-facing setting:
the learned field is evaluated on the same full grid and with the same
operator used by FGMRES.

% ============================================================
\section{Discrete Helmholtz systems}
\label{sec:setup-discrete-systems}
% ============================================================

All experiments solve linear systems of the form
\begin{equation}
    A_\omega u_\omega = b,
    \label{eq:setup-linear-system}
\end{equation}
where $\omega$ is the angular frequency, $u_\omega$ is the complex-valued
wavefield, and $b$ is the complex-valued source. The residual convention is
\begin{equation}
    r(x)=b-A_\omega x.
    \label{eq:setup-residual}
\end{equation}
This convention is used for training diagnostics, warm-start evaluation, and
FGMRES convergence plots. It is important because the central question is not
only whether a neural prediction resembles the high-frequency solution, but
whether it is compatible with the high-frequency discrete equation.

% ------------------------------------------------------------
\subsection{One-dimensional Dirichlet problem}
\label{subsec:setup-1d-dirichlet}
% ------------------------------------------------------------

The clean diagnostic problem is the homogeneous one-dimensional Dirichlet
Helmholtz system with $N=512$ unknowns and no absorbing layer. The tested
frequency pair is
\begin{equation}
    \omega_L=16,
    \qquad
    \omega_H=32.
    \label{eq:setup-1d-dirichlet-pair}
\end{equation}
With the sign convention used throughout the one-dimensional spectral
analysis,
\begin{equation}
    A_\omega=-D_{xx}-\omega^2 I.
    \label{eq:setup-1d-dirichlet-operator}
\end{equation}
The eigenpairs are analytical:
\begin{equation}
    \lambda_k(\omega)
    =
    \frac{4}{h^2}
    \sin^2\!\left(\frac{\pi k}{2(N+1)}\right)
    -
    \omega^2,
    \qquad
    v_k(j)
    =
    \sqrt{\frac{2}{N+1}}
    \sin\!\left(\frac{j k \pi}{N+1}\right),
    \label{eq:setup-dirichlet-eigenpairs}
\end{equation}
for $j,k=1,\dots,N$. The eigenvectors are Euclidean-normalised. Therefore, if
$u_\star$ is the high-frequency solution, $x_0$ is an initial guess, and
$e_0=u_\star-x_0$, then
\begin{equation}
    r_0=b-A_Hx_0=A_He_0,
    \qquad
    c_k(r_0)=\lambda_k c_k(e_0).
    \label{eq:setup-modal-residual-map}
\end{equation}
This identity is the reason for including the Dirichlet problem. It turns the
field-error/residual-error distinction from a qualitative warning into an
exact mode-by-mode diagnostic.

% ------------------------------------------------------------
\subsection{One-dimensional PML problem}
\label{subsec:setup-1d-pml}
% ------------------------------------------------------------

The one-dimensional PML problem uses the same frequency pair, $16\to32$, but
replaces the simple Dirichlet boundary treatment by an absorbing layer. This
setting is deliberately intermediate. It is still small enough for careful
inspection and repeated solver runs, but the simple symmetric sine-basis
analysis no longer applies in the same way because the PML operator is
complex-valued and generally non-normal.

The role of this experiment is to test operator compatibility in the absorbing
layer. A predicted field can be good in the physical region while producing a
large full-grid residual if the PML values are inconsistent with the discrete
stencil and stretching coefficients. For this reason, the PML experiment
reports full-grid field errors, true residuals, preconditioned residuals, and
FGMRES curves, but does not present the PML problem as an exact modal
decomposition.

% ------------------------------------------------------------
\subsection{Two-dimensional FD/PML problem}
\label{subsec:setup-2d-fd-pml}
% ------------------------------------------------------------

The solver-facing experiments use a two-dimensional finite-difference
Helmholtz operator on a full $512\times512$ computational grid. A PML of width
$112$ grid points surrounds the physical domain on each side, leaving a
$288\times288$ interior. The tested frequency pairs are
\begin{equation}
    16\to32,
    \qquad
    32\to64,
    \qquad
    64\to128.
    \label{eq:setup-2d-pairs}
\end{equation}

The PML is treated as part of the algebraic system, not as a removable visual
border. All solver-facing residuals are computed on the full $512\times512$
grid using the same FD/PML operator used by FGMRES. Interior metrics are
reported only as additional diagnostics. The primary question is whether the
predicted full-grid field gives a useful state for the original
high-frequency equation.

\begin{table}[htbp]
\centering
\caption{Discrete systems used in the experiments. The Dirichlet problem gives
an exact spectral diagnostic, the 1D PML problem isolates absorbing-layer
compatibility, and the 2D FD/PML problem is the full solver-facing setting.}
\label{tab:setup-systems}
\begin{tabular}{lcccl}
\toprule
Setting & Grid & PML width & Interior & Frequency pairs \\
\midrule
1D Dirichlet & $512$ & none & $512$ & $16\to32$ \\
1D PML & problem-dependent & included & physical subdomain & $16\to32$ \\
2D FD/PML & $512\times512$ & $112$ per side & $288\times288$ &
$16\to32$, $32\to64$, $64\to128$ \\
\bottomrule
\end{tabular}
\end{table}

\begin{figure}[htbp]
    \centering
    \chsetupincludegraphics{0.88\linewidth}{figures/ch6/discrete_systems_overview.png}
    \caption{Recommended overview figure for the three computational systems:
    1D Dirichlet, 1D PML, and 2D FD/PML. The figure should show which region
    belongs to the physical domain, which region belongs to the PML when
    present, and which operator is used to compute the residual.}
    \label{fig:setup-discrete-systems-overview}
\end{figure}

% ============================================================
\section{Datasets and source conventions}
\label{sec:setup-datasets}
% ============================================================

Each supervised frequency-transfer sample contains a source, a
low-frequency solution, and the corresponding high-frequency solution. Complex
quantities are stored by real and imaginary channels. The source convention is
therefore part of the data specification: both source channels are retained so
that the learned model and the residual evaluation refer to the same complex
linear system.

The two-dimensional FD/PML datasets contain $9600$ full-grid samples for each
frequency pair. Each sample uses a complex Gaussian source and stores the
low- and high-frequency solutions on the full $512\times512$ grid. The
imaginary source channel is stored. Dataset audits were performed before the
solver-facing evaluations, so the data used for the neural models and the
data used for residual evaluation refer to the same complex-valued FD/PML
convention.

The full-grid convention is a deliberate correction to earlier
interior-focused workflows. A model trained only on the physical interior can
appear accurate in plots while leaving the absorbing layer incompatible with
the operator. The modern datasets therefore keep the PML part of the
wavefield in the training target and in the residual evaluation.

\begin{table}[htbp]
\centering
\caption{Two-dimensional FD/PML dataset specification. Each transfer pair has
$9600$ full-grid samples with complex Gaussian sources.}
\label{tab:setup-2d-datasets}
\begin{tabular}{lccc}
\toprule
Pair & Samples & Stored source channels & Stored solution region \\
\midrule
$16\to32$ & $9600$ & real and imaginary & full $512\times512$ grid \\
$32\to64$ & $9600$ & real and imaginary & full $512\times512$ grid \\
$64\to128$ & $9600$ & real and imaginary & full $512\times512$ grid \\
\bottomrule
\end{tabular}
\end{table}

The one-dimensional datasets are used for diagnostic rather than production
scale. The Dirichlet data support exact modal analysis, spectral filtering,
and residual gating. The 1D PML data support the comparison between older
PML-handling conventions and the corrected full-flux PML convention used in
the Results chapter.

% ============================================================
\section{Neural transfer models}
\label{sec:setup-neural-models}
% ============================================================

The learned models are supervised frequency-transfer maps. In the warm-start
experiments, the model predicts a high-frequency field from a lower-frequency
field and the prediction is used as the initial guess for FGMRES. In the
one-dimensional diagnostic experiments, the same learned information is also
filtered, gated, or inserted into learned preconditioning variants in order to
test whether it can be made residual-compatible.

Complex fields are represented in real arithmetic by separate real and
imaginary channels. The basic field loss is the complex relative $\ell^2$
error
\begin{equation}
    \mathcal{L}_{\mathrm{field}}
    =
    \frac{\|\hat{u}_H-u_H\|_2}{\|u_H\|_2}.
    \label{eq:setup-field-loss}
\end{equation}
This objective measures field prediction quality. It does not by itself
guarantee a small residual under $A_H$.

% ------------------------------------------------------------
\subsection{One-dimensional models}
\label{subsec:setup-1d-models}
% ------------------------------------------------------------

The one-dimensional Dirichlet model is a U-Net trained to map $u_{16}$ to
$u_{32}$. The final model uses the real and imaginary parts of the
low-frequency field together with real and imaginary right-hand-side channels.
The target is the complex high-frequency solution. Its best validation field
loss is approximately $5.37\times10^{-3}$, so it is a strong field predictor
in the controlled Dirichlet setting.

The Results chapter uses this model in several diagnostic roles. First, it is
used raw as a warm start. Second, its modal content is filtered or gated in
the known Dirichlet eigenbasis. Additional residual-correction experiments are
treated as diagnostic extensions: they separate residual restriction,
low-frequency solving, and high-frequency correction prolongation, but they
are not the main solver claim of the thesis. Their role is to test whether the
field-trained information can be made safe for the residual.

The one-dimensional PML comparison uses learned starts that differ mainly in
how they handle the absorbing layer. The older green-function/PML variants are
compared with a corrected flux-full model. The key distinction is whether the
PML output is kept, zeroed, or trained as part of the full
operator-compatible field.

% ------------------------------------------------------------
\subsection{Two-dimensional full-grid model}
\label{subsec:setup-2d-model}
% ------------------------------------------------------------

The main two-dimensional model is a full-grid flux-trained
\texttt{TransferUNet} with \texttt{base\_ch=32} and five levels. The input is
the real and imaginary parts of the low-frequency field. The target is the
real and imaginary parts of the high-frequency field on the full
$512\times512$ FD/PML grid. The loss is again the complex relative $\ell^2$
field loss, evaluated on the full grid. Older interior-loss models are kept as
baselines because they expose the PML failure mode: they can predict the
physical interior well while producing an absorbing-layer output that is
incompatible with the full operator.

\begin{table}[htbp]
\centering
\caption{Neural transfer models used in the experiments. The models are field
predictors by training objective; solver usefulness is tested separately by
residuals and FGMRES curves.}
\label{tab:setup-models}
\begin{tabular}{llll}
\toprule
Setting & Model & Input channels & Training target \\
\midrule
1D Dirichlet & U-Net & $u_L$ and RHS, real/imag & full-grid $u_H$ \\
1D PML & PML transfer variants & $u_L$-based channels & full/PML-dependent $u_H$ \\
2D FD/PML & \texttt{TransferUNet}, \texttt{base\_ch=32}, 5 levels &
$u_L$, real/imag & full-grid $u_H$ \\
\bottomrule
\end{tabular}
\end{table}

\begin{figure}[htbp]
    \centering
    \chsetupincludegraphics{0.90\linewidth}{figures/ch6/model_and_data_pipeline.png}
    \caption{Model and data pipeline used for learned frequency transfer. Complex sources generate paired low- and high-frequency solutions; the U-Net is trained as a field predictor using real and imaginary channels; and the resulting warm start is evaluated with solver-native residuals and FGMRES behaviour.}
    \label{fig:setup-model-data-pipeline}
\end{figure}

% ============================================================
\section{Warm starts and learned correction tests}
\label{sec:setup-warm-starts-corrections}
% ============================================================

The simplest learned solver use is a warm start:
\begin{equation}
    x_0=\hat{u}_H=T_{\uparrow,\theta}^{\mathrm{sol}}u_L.
    \label{eq:setup-warm-start}
\end{equation}
The cold-start baseline is $x_0=0$. For a warm start, the initial true
relative residual is
\begin{equation}
    r_0^{\mathrm{true}}
    =
    \frac{\|b-A_Hx_0\|_2}{\|b\|_2}.
    \label{eq:setup-true-initial-residual}
\end{equation}
This is the headline algebraic diagnostic for the initial state.

The one-dimensional Dirichlet experiments also test residual-aware selection.
Because the eigenbasis is known, a learned modal component can be accepted
only when it reduces the corresponding residual contribution relative to the
baseline. This residual gate is used as an explanatory device. It demonstrates
that a field-trained model may contain useful solver-facing components even
when the raw prediction is not residual-safe.

For learned correction tests, the mathematical role changes. A
preconditioner receives a residual and returns a correction. The relevant
ideal is therefore
\begin{equation}
    A_H z_H \approx r_H,
    \label{eq:setup-correction-task}
\end{equation}
not direct prediction of a complete high-frequency solution. The diagnostic
experiments distinguish three objects:
\begin{align}
    T_{\uparrow}^{\mathrm{sol}} &: u_L \mapsto u_H,
    &&\text{solution transfer}, \label{eq:setup-object-solution}\\
    R &: r_H \mapsto r_L \text{ or } e_L,
    &&\text{residual restriction}, \label{eq:setup-object-restriction}\\
    P &: e_L \mapsto e_H,
    &&\text{correction prolongation}. \label{eq:setup-object-prolongation}
\end{align}
These roles match the conceptual separation introduced in
Chapter~\ref{ch:frequency-transfer}. In the Results chapter this distinction
is used mainly to interpret the warm-start and gating results; extended
residual-correction tests are best reported as appendix material unless they
produce a clear iteration-count improvement.

% ============================================================
\section{CSL-FGMRES solver protocol}
\label{sec:setup-solver-protocol}
% ============================================================

The solver-facing experiments use FGMRES with a complex shifted Laplace (CSL)
preconditioner. For the high-frequency operator $A$, the shifted operator has
the same discrete Laplacian and a complex-shifted mass term, schematically
\begin{equation}
    M_\beta
    =
    -\Delta_h-(1+i\beta)\omega^2 I,
    \label{eq:setup-csl-operator}
\end{equation}
with the same boundary or PML convention as the corresponding Helmholtz
operator. The preconditioner applies $M_\beta^{-1}$ through an exact sparse
LU solve of the shifted operator. In the reported main results, the damping is
fixed at
\begin{equation}
    \beta=0.3.
    \label{eq:setup-csl-beta}
\end{equation}
Earlier $\beta=0.1$ experiments are not used for the main comparison because
mixing damping values would obscure whether a difference comes from the
learned start or from the classical preconditioner.

For a fixed frequency, grid, PML convention, and damping value, the sparse
operator and CSL factorisation are assembled once and reused across
right-hand sides whenever possible. This keeps the comparison focused on the
initial state and Krylov behaviour rather than repeated setup cost. In
warm-start comparisons, the operator, right-hand side, preconditioner,
tolerance, and iteration budget are held fixed; only the initial guess
changes.

The two-dimensional solver evaluation uses $10$ test samples per frequency
pair and a fixed budget of $40$ FGMRES iterations. At $\beta=0.3$, the
reported 2D runs do not reach the target tolerance within this budget.
Therefore, the finite-budget final residual is the meaningful 2D comparison,
not the nominal iteration count to convergence.

\begin{table}[htbp]
\centering
\caption{Solver-evaluation protocol for the main comparisons in
Chapter~\ref{chapter:results}.}
\label{tab:setup-solver-protocol}
\begin{tabular}{ll}
\toprule
Quantity & Protocol \\
\midrule
Krylov method & FGMRES \\
Cold start & $x_0=0$ \\
Neural warm start & $x_0=\hat{u}_H$ \\
Primary initial metric & $\|b-A_Hx_0\|_2/\|b\|_2$ \\
Classical preconditioner & CSL \\
Main CSL damping & $\beta=0.3$ \\
CSL application & exact sparse LU solve of shifted operator \\
2D test samples & $10$ per frequency pair \\
2D iteration budget & $40$ iterations \\
2D main comparison & finite-budget final true residual \\
\bottomrule
\end{tabular}
\end{table}

% ------------------------------------------------------------
\subsection{Left and right preconditioning}
\label{subsec:setup-left-right-preconditioning}
% ------------------------------------------------------------

The Results chapter includes an explicit left-versus-right preconditioning
comparison for the one-dimensional Dirichlet problem. The distinction matters
because the residual minimised by GMRES depends on how the preconditioner is
inserted. Left preconditioning solves
\begin{equation}
    M^{-1}Ax=M^{-1}b,
    \label{eq:setup-left-preconditioning}
\end{equation}
whereas right preconditioning solves
\begin{equation}
    AM^{-1}y=b,
    \qquad
    x=M^{-1}y.
    \label{eq:setup-right-preconditioning}
\end{equation}

The primary thesis metric is always the true residual of the original
Helmholtz system,
\begin{equation}
    \frac{\|b-Ax_k\|_2}{\|b\|_2}.
    \label{eq:setup-true-residual}
\end{equation}
The preconditioned residual
\begin{equation}
    \frac{\|M^{-1}(b-Ax_k)\|_2}{\|M^{-1}b\|_2}
    \label{eq:setup-preconditioned-residual}
\end{equation}
is also reported when it explains what the preconditioner presents to the
Krylov process. This separation prevents a left-preconditioned curve from
being mistaken for the true residual of the original Helmholtz equation.

\begin{figure}[htbp]
    \centering
    \chsetupincludegraphics{0.86\linewidth}{figures/ch6/residual_norms_schematic.png}
    \caption{Recommended residual-norm schematic. The figure should contrast
    the true residual $\|b-Ax_k\|_2/\|b\|_2$ with the preconditioned residual
    $\|M^{-1}(b-Ax_k)\|_2/\|M^{-1}b\|_2$, and indicate why left and right
    preconditioning can make these curves tell different stories.}
    \label{fig:setup-residual-norms-schematic}
\end{figure}

% ============================================================
\section{Diagnostics and metrics}
\label{sec:setup-metrics}
% ============================================================

The thesis reports several diagnostic quantities because no single scalar
captures the behaviour of a learned Helmholtz solver component.

\paragraph{Field error.}
The relative field error is
\begin{equation}
    E_{\mathrm{field}}
    =
    \frac{\|\hat{u}_H-u_H\|_2}{\|u_H\|_2}.
    \label{eq:setup-field-error}
\end{equation}
In 2D this may be reported on both the full grid and the physical interior.
The full-grid value is the primary quantity for the modern FD/PML models,
because the solver applies the full-grid operator.

\paragraph{Initial true residual.}
The initial true residual is
\begin{equation}
    E_{\mathrm{res},0}
    =
    \frac{\|b-A_Hx_0\|_2}{\|b\|_2}.
    \label{eq:setup-initial-true-residual}
\end{equation}
This is the first solver-facing test of a learned warm start. A prediction can
have small field error and still have a worse initial true residual than the
cold start.

\paragraph{Preconditioned residual.}
When CSL is active, the preconditioned residual
\begin{equation}
    E_{\mathrm{prec},0}
    =
    \frac{\|M^{-1}(b-A_Hx_0)\|_2}{\|M^{-1}b\|_2}
    \label{eq:setup-initial-preconditioned-residual}
\end{equation}
is reported as an explanatory diagnostic. It is especially useful when a raw
learned start has a large true residual but a much smaller residual after CSL
preconditioning.

\paragraph{FGMRES curves and finite-budget residuals.}
Convergence curves are reported because different starts can enter FGMRES
with very different residual levels and then converge with similar later
slopes. For the 2D FD/PML experiments at $\beta=0.3$, no method reaches the
target tolerance within the fixed $40$-iteration budget, so the final residual
after that budget is the main solver-facing comparison.

\paragraph{Modal diagnostics.}
In the 1D Dirichlet problem, the sine eigenbasis gives exact modal
coefficients. This makes it possible to identify which field-error components
are amplified into residual components and to test residual gating mode by
mode.

\paragraph{PML diagnostics.}
In PML settings, the absorbing layer is diagnosed by full-grid residuals and,
when useful, by ratios that compare PML-region field magnitude with full-grid
or interior behaviour. These diagnostics are included because zeroing or
mishandling the PML can improve the appearance of the physical-domain plot
while damaging the operator residual.

\begin{figure}[htbp]
    \centering
    \chsetupincludegraphics{0.92\linewidth}{figures/ch6/metric_dashboard_example.png}
    \caption{Recommended metric-dashboard figure. For one representative
    sample, show field error, true initial residual, preconditioned residual,
    and the first part of the FGMRES curve. This would make clear why the
    thesis does not judge learned transfer by field loss alone.}
    \label{fig:setup-metric-dashboard}
\end{figure}

% ============================================================
\section{Computational workflow and artifacts}
\label{sec:setup-workflow-artifacts}
% ============================================================

The computational workflow has four stages: dataset generation,
neural-network training, solver evaluation, and diagnostic analysis. Dataset
generation and solver evaluation are dominated by sparse Helmholtz solves.
Neural-network training is dominated by GPU throughput. One-dimensional
spectral diagnostics are small enough to support dense checks or analytical
eigenbasis calculations.

\begin{table}[htbp]
\centering
\caption{Computational stages used in the experiments. The final evaluation is
solver-native: field loss is recorded, but residuals and FGMRES behaviour
decide whether a variant is useful.}
\label{tab:setup-computational-stages}
\begin{tabular}{lll}
\toprule
Stage & Main bottleneck & Primary outputs \\
\midrule
Dataset generation & Sparse solves and I/O & fields, sources, metadata \\
Neural-network training & GPU throughput & checkpoints, validation losses \\
Solver evaluation & Sparse solves and FGMRES & residuals, iterations, curves \\
Diagnostic analysis & spectra, plots, summaries & modal/PML diagnostics and figures \\
\bottomrule
\end{tabular}
\end{table}

Experiment outputs are kept in named directories that record the frequency
pair, boundary/PML convention, training objective, or diagnostic family.
Corrected pipelines are stored separately from exploratory pipelines rather
than overwriting them. This matters because several conventions can produce
similar-looking wavefield plots while representing different algebraic
systems. Green-function data, interior-focused PML data, row-scaled PML data,
and corrected full-grid FD/PML data should therefore not be mixed without
explicit labels.

For thesis-quality runs, the reproducibility record consists of the launch
script or sweep specification, numerical metadata, trained checkpoint, summary
table, residual curves, and figure-generation script. Training checkpoints
such as \texttt{best.pt} and \texttt{last.pt} are retained when applicable, so
that a model can be evaluated or resumed without reconstructing the run from
terminal history. Derived quantities such as residuals, modal coefficients,
and plots are treated as regenerable analysis artifacts rather than primary
dataset fields.

The main one-dimensional artifacts used by the Results chapter are stored
under the corrected flux-pipeline output tree,
\begin{center}
\texttt{experiments/claude/eigenvalue\_1d/corrected\_flux\_pipeline/}.
\end{center}
The final chapter figures are copied from the relevant analysis folders into
\texttt{figures/ch7/}. The important reproducibility point is the naming
convention: the figure names identify whether the plot is Dirichlet or PML,
whether it is a true or preconditioned residual, and which frequency pair and
CSL damping value were used.

The main presentation plots for two-dimensional experiments are stored under
\begin{center}
\texttt{experiments/2d/presentation\_plots/}.
\end{center}
The full-grid training curves use the \texttt{flux\_full\_N9600} outputs, and
the solver-facing curves use the corresponding full-grid solver-evaluation
outputs.
The figure captions in Chapter~\ref{chapter:results} record the copied figure
name used for each reported claim, while the experiment directories retain the
regenerable scripts, summaries, and raw curves.

\begin{table}[htbp]
\centering
\caption{Main computational risks and safeguards in the modern experiments.}
\label{tab:setup-computational-safeguards}
\begin{tabular}{lll}
\toprule
Risk & Consequence & Safeguard \\
\midrule
Wrong discretisation convention & misleading residuals or spectra &
validate in 1D first \\
Field loss improves only visually & no Krylov benefit &
evaluate true residuals and FGMRES curves \\
PML incompatibility & inflated full-grid residual &
train and evaluate full-grid FD/PML fields \\
Mixed CSL damping values & ambiguous solver comparison &
use $\beta=0.3$ for main results \\
Untraceable sweeps & irreproducible comparisons &
use named output folders and summaries \\
Repeated setup costs & distorted timing interpretation &
reuse operators and factorizations \\
\bottomrule
\end{tabular}
\end{table}

The hardware organisation follows from these stages but is not itself a
scientific variable. Local runs are used for syntax checks, small validation
tests, and plotting. Interactive GPU servers are used for training and rapid
solver debugging. Cluster resources are used for long sweeps and repeated
benchmarks. Across all environments, the rule is to use the smallest
computation that answers the current question reliably and reserve large runs
for designs that have already passed small-scale validation.

% ============================================================
\section{Summary}
\label{sec:setup-summary}
% ============================================================

This chapter defined the computational setup used to test learned frequency
transfer as a solver component. The one-dimensional Dirichlet problem provides
an exact spectral diagnostic for residual amplification. The one-dimensional
PML problem tests absorbing-layer compatibility without the full cost and
complexity of the two-dimensional setting. The two-dimensional FD/PML problem
is the main solver-facing evaluation, with full-grid training targets and
full-grid residuals.

The setup separates three questions that are easy to conflate. First, can a
neural model predict a high-frequency field from a low-frequency field?
Second, does the predicted field have a small residual under the
high-frequency operator? Third, does the resulting initial state or learned
correction reduce FGMRES work? The Results chapter answers these questions by
reporting field errors, true residuals, preconditioned residuals, modal/PML
diagnostics, and finite-budget FGMRES behaviour.

% ============================================================
% Recommended visual plan for this chapter
% ============================================================
%
% Essential figures:
% 1. figures/ch6/discrete_systems_overview.png
%    Three panels: 1D Dirichlet, 1D PML, 2D FD/PML.
%    Purpose: show why the experiments form a ladder from exact diagnostics
%    to full solver-facing evaluation.
%
% 2. figures/ch6/model_and_data_pipeline.png
%    Source -> low/high solves -> real/imag channels -> U-Net -> residual
%    evaluation.
%    Purpose: show that the model is trained as a field predictor but judged
%    by the operator residual.
%
% 3. figures/ch6/residual_norms_schematic.png
%    True residual versus preconditioned residual, with left/right
%    preconditioning indicated.
%    Purpose: prepare the reader for the left/right comparison in the Results
%    chapter.
%
% Optional figure:
% 4. figures/ch6/metric_dashboard_example.png
%    One representative sample with field error, true residual,
%    preconditioned residual, and early FGMRES curve.
%    Purpose: improve readability by making the metric hierarchy concrete.
